
We henceforth assume that \(S\) is the spectrum of a field \(k\), and
that \(G\), \(V\), and \(W\) are of finite type over \(k\). Let
\(\mathfrak n\) be the Lie algebra of \(N\). Then
\(\dim N\leq\operatorname{rank}_k\mathfrak n\), with equality if and
only if \(N\) is smooth over \(k\) (Exposé VI). Let \(a\in W(k)\), and
consider the subscheme \(M_a\) of \(G\) defined above; it contains
\(N\) and is isomorphic to \(\varphi^{-1}(a)\). We denote by
\(\mathfrak m_a\) its Zariski tangent space at the identity element
\(e\) of \(N\). Thus\footnote{The French source omits
\(\operatorname{rank}_k\) before \(\mathfrak m_a\) in \SGARefLink{sga3:xiii:equation:1}{(1)}; it is
restored here, as required by the surrounding text and \SGARefLink{sga3:xiii:lemma:1:1}{Lemma 1.1}.}
\SGARefTarget{sga3:xiii:equation:1}
\[
\tag{1}
  \mathfrak n\subset\mathfrak m_a,\qquad
  \dim N\leq\operatorname{rank}_k\mathfrak n
  \leq\operatorname{rank}_k\mathfrak m_a.
\]

\medskip
\noindent\SGARefTarget{sga3:xiii:lemma:1:1}\textbf{Lemma 1.1.} ---
\textit{With the preceding notation:}

\begin{enumerate}
\item[\textit{a)}] \SGARefTarget{sga3:xiii:lemma:1:1:item:a} Consider the following conditions:
  \begin{enumerate}
  \item[(i)] \SGARefTarget{sga3:xiii:lemma:1:1:item:i}
    \(\mathfrak n=\mathfrak m_a\), and \(N\) is smooth over \(k\).
  \item[(i bis)] \SGARefTarget{sga3:xiii:lemma:1:1:item:i-bis}
    \(\dim N=\operatorname{rank}_k\mathfrak m_a\).
  \item[(ii)] \SGARefTarget{sga3:xiii:lemma:1:1:item:ii}
    The morphism \(\psi:X\to V\) is unramified at
    \((\bar e,a)\).
  \item[(iii)] \SGARefTarget{sga3:xiii:lemma:1:1:item:iii}
    \(M_a\) and \(N\) coincide in a neighborhood of \(e\).
  \end{enumerate}
  Then
  \[
    \SGARefLink{sga3:xiii:lemma:1:1:item:i}{\mathrm{(i)}}\Longleftrightarrow\SGARefLink{sga3:xiii:lemma:1:1:item:i-bis}{\mathrm{(i\ bis)}}
    \Longrightarrow\SGARefLink{sga3:xiii:lemma:1:1:item:ii}{\mathrm{(ii)}}
    \Longleftrightarrow\SGARefLink{sga3:xiii:lemma:1:1:item:iii}{\mathrm{(iii)}}.
  \]
\item[\textit{b)}] \SGARefTarget{sga3:xiii:lemma:1:1:item:b}
  Suppose that \(\varphi:G\times_S W\to V\) is smooth at \((e,a)\).
  Then \(M_a\) is smooth over \(k\) at \(e\), and \(\psi\) is smooth
  at \((\bar e,a)\).
\end{enumerate}

\noindent\textit{Proof.}
\SGARefTarget{sga3:xiii:lemma:1:1:proof:item:a}
\textit{a)} The equivalence of \SGARefLink{sga3:xiii:lemma:1:1:item:i}{(i)} and \SGARefLink{sga3:xiii:lemma:1:1:item:i-bis}{(i bis)} follows immediately
from \SGARefLink{sga3:xiii:equation:1}{(1)} and from the fact noted above that \(N\) is smooth over \(k\)
if and only if
\(\dim N=\operatorname{rank}_k\mathfrak n\). Consider the inclusion
morphism \(N\to M_a\). It is well known that if \(N\) is smooth over \(k\) at \(e\) and
the tangent map at \(e\) is surjective, then \(N\to M_a\) is smooth
at \(e\); since it is an immersion, it is therefore an isomorphism at
\(e\). This proves that \SGARefLink{sga3:xiii:lemma:1:1:item:i}{(i)} implies \SGARefLink{sga3:xiii:lemma:1:1:item:iii}{(iii)}.

To prove the equivalence of \SGARefLink{sga3:xiii:lemma:1:1:item:ii}{(ii)} and \SGARefLink{sga3:xiii:lemma:1:1:item:iii}{(iii)}, let
\(X_a=\psi^{-1}(a)\), as above. From the isomorphisms
\[
  M_a\simeq\varphi^{-1}(a)\simeq q^{-1}(X_a)
\]
we obtain a morphism \(p_a:M_a\to X_a\) making \(M_a\) a principal
homogeneous bundle with group \(N_{X_a}\). Consider \SGARefLink{sga3:xiii:lemma:1:1:diagram:01}{the diagram}
\SGARefTarget{sga3:xiii:lemma:1:1:diagram:01}
\[
\begin{tikzcd}[column sep=large,row sep=large]
M_a \arrow[d,"p_a"'] &
N \arrow[l,"j_a"'] \arrow[d] \\
X_a \arrow[r,"j'_a"] &
S=\operatorname{Spec}(k).
\end{tikzcd}
\]
Here \(\operatorname{Spec}(k)\to X_a\)\footnote{\SGARefLink{sga3:xiii:lemma:1:1:diagram:01}{The printed diagram}
directs \(j'_a:X_a\to S\), whereas the immediately following prose
defines the point morphism \(S\to X_a\) and calls it \(j'_a\). Both
\SGARefLink{sga3:xiii:lemma:1:1:diagram:01}{the printed diagram direction} and the verbal definition are retained
here. The source also prints \(X_n\) here and \(X_{a,b}\) in the two
local rings below. They are corrected to \(X_a\) and
\(X_{a,\bar a}\), in accordance with the fiber notation and the point
\(\bar a\).} is defined by the point
\(\bar a=(\bar e,a)\) of \(X\), and \(j_a:N\to M_a\) is the canonical
immersion. To say that \(\psi\) is unramified at \((\bar e,a)\) means
that \(j'_a\) is an open immersion, or equivalently that it induces an
isomorphism
\SGARefTarget{sga3:xiii:lemma:1:1:display:residue-field-map}
\[
  S\xrightarrow{\sim}\operatorname{Spec}
  \bigl(\mathcal O_{X_a,\bar a}\bigr).
\]
Since \(p_a\) is flat, this is equivalent to saying that the morphism
obtained from \SGARefLink{sga3:xiii:lemma:1:1:display:residue-field-map}{the preceding one} by the base change
\[
  \operatorname{Spec}\bigl(\mathcal O_{M_a,e}\bigr)
  \longrightarrow
  \operatorname{Spec}\bigl(\mathcal O_{X_a,\bar a}\bigr)
\]
is an isomorphism. The resulting morphism is precisely
\[
  \operatorname{Spec}\bigl(\mathcal O_{N,e}\bigr)
  \longrightarrow
  \operatorname{Spec}\bigl(\mathcal O_{M_a,e}\bigr),
\]
which proves the equivalence of \SGARefLink{sga3:xiii:lemma:1:1:item:ii}{(ii)} and \SGARefLink{sga3:xiii:lemma:1:1:item:iii}{(iii)}.

Notice moreover that the proof shows that \SGARefLink{sga3:xiii:lemma:1:1:item:ii}{(ii)} and \SGARefLink{sga3:xiii:lemma:1:1:item:iii}{(iii)} imply \SGARefLink{sga3:xiii:lemma:1:1:item:iii-bis}{the
following condition}, apparently stronger than \SGARefLink{sga3:xiii:lemma:1:1:item:iii}{(iii)}:
\begin{enumerate}
\item[(iii bis)] \SGARefTarget{sga3:xiii:lemma:1:1:item:iii-bis}
  \(N\) is an open and closed subscheme of \(M_a\).
\end{enumerate}

\SGARefTarget{sga3:xiii:lemma:1:1:proof:item:b}
\textit{b)} The first assertion follows from the isomorphism
\(M_a\simeq\varphi^{-1}(a)\), and the second from the fact that \(q\)
is flat and \(q(e,a)=(\bar e,a)\).

\section{Density theorem and theory of the regular points of \(G\)}
\SGARefTarget{sga3:xiii:section:2}

We apply the constructions and notation of \SGARefLink{sga3:xiii:section:1}{the preceding section} in
the following situation. Let \(G\) be a smooth connected algebraic
group over \(k\); take \(V=G\), with \(G\) acting on \(V\) by inner
automorphisms; and let \(W=H\), where \(H\) is a smooth connected
algebraic subgroup of \(G\). We denote the Lie algebras of \(G\) and
\(H\) by \(\mathfrak g\) and \(\mathfrak h\), the normalizer of \(H\)
in \(G\) by \(N\), and its Lie algebra by \(\mathfrak n\). For
\(a\in G(k)\), we again denote by \(M_a\), as in \SGARefLink{sga3:xiii:section:1}{Section 1}, the
inverse of the transporter of \(a\) into \(H\). Thus, if
\(a\in H(k)\), then \(N\subset M_a\); in that case
\(\mathfrak m_a\) denotes the Zariski tangent space of \(M_a\) at the
identity element \(e\) of \(G\). Note that
\[
  \mathfrak h\subset\mathfrak n\subset\mathfrak m_a
  \subset\mathfrak g
  \qquad\text{for }a\in H(k).
\]
We shall use \SGARefLink{sga3:xiii:lemma:2:0}{the following lemma}.

\medskip
\noindent\SGARefTarget{sga3:xiii:lemma:2:0}\textbf{Lemma 2.0.} ---
\textit{The equality \(H=N^0\) holds if and only if}
\SGARefTarget{sga3:xiii:lemma:2:0:display:invariants-zero}
\[
  (\mathfrak g/\mathfrak h)^H=0,
\]
\textit{where the left-hand side denotes the subspace of invariants
for the action of \(H\) induced by the adjoint representation. When
\SGARefLink{sga3:xiii:lemma:2:0:display:invariants-zero}{this condition} holds, \(N\) is smooth and \(\dim X=\dim G\). In all
cases, \(\dim X\leq\dim G\), and equality holds if and only if \(H\)
has finite index in \(N\).}

Indeed, by \SGARefLink{sga3:ii:theorem:5-2-3:item:i}{II~5.2.3(i)}, \(\mathfrak n\) is the inverse image of
\((\mathfrak g/\mathfrak h)^H\) under
\(\mathfrak g\to\mathfrak g/\mathfrak h\). Hence
\((\mathfrak g/\mathfrak h)^H=0\) is equivalent to
\(\mathfrak h=\mathfrak n\), which in turn is equivalent to
\(H=N^0\), since \(H\) is a smooth connected algebraic subgroup of
the algebraic group \(N\) (cf. VI.2). This also shows that \(N\) is
smooth. Furthermore,
\[
  \dim X=(\dim G-\dim N)+\dim H
        =\dim G-(\dim N-\dim H),
\]
and therefore
\[
  \dim X\leq\dim G,
\]
with equality if and only if \(\dim H=\dim N\), that is, if and only
if \(H\) has finite index in \(N\). This holds in particular when
\(H=N^0\), and proves the lemma.

\medskip
\noindent\SGARefTarget{sga3:xiii:theorem:2:1}\textbf{Theorem 2.1.} ---
\textit{Let \(G\) be a smooth connected algebraic group over an
algebraically closed field \(k\), let \(H\) be a smooth connected
algebraic subgroup, and let \(N\) be its normalizer. Write
\(\mathfrak g,\mathfrak h,\mathfrak n\) for their Lie algebras, let
\(X=G\times^N H\) be the scheme introduced in \SGARefLink{sga3:xiii:section:1}{Section 1} (fibered over
\(G/N\) with typical fiber \(H\)), and let \(\psi:X\to G\) be the
canonical morphism. Its image is also the image of
\(\varphi:G\times H\to G\), where}
\[
  \varphi(g,h)=\operatorname{int}(g)h=ghg^{-1}.
\]
\textit{The following conditions are equivalent:}

\begin{enumerate}
\item[(i)] \SGARefTarget{sga3:xiii:theorem:2:1:item:i}
  \(H\) contains a Cartan subgroup \(C\) of \(G\) (XII~1).
\item[(i bis)] \SGARefTarget{sga3:xiii:theorem:2:1:item:i-bis}
  \(H\) has the same reductive rank and the same nilpotent rank
  (XII~1) as \(G\).
\item[(ii)] \SGARefTarget{sga3:xiii:theorem:2:1:item:ii}
  \(H\) contains a maximal torus \(T\) of \(G\), and
  \((\mathfrak g/\mathfrak h)^T=0\).
\item[(iii)] \SGARefTarget{sga3:xiii:theorem:2:1:item:iii}
  The set of conjugates of \(H\) containing a given maximal torus is
  finite and nonempty, and \(H\) has finite index in its normalizer.
\item[(iv)] \SGARefTarget{sga3:xiii:theorem:2:1:item:iv}
  There exists \(a\in H(k)\) that belongs to only finitely many
  conjugates of \(H\) (or merely such that \(\psi^{-1}(a)\) has an
  isolated point), and \(H\) has finite index in its normalizer.
\item[(iv bis)] \SGARefTarget{sga3:xiii:theorem:2:1:item:iv-bis}
  The morphism \(\psi:X\to G\) is generically quasi-finite (that is,
  there is a dense open subset of \(X\) on which \(\psi\) is
  quasi-finite), and \(H\) has finite index in its normalizer.
\item[(v)] \SGARefTarget{sga3:xiii:theorem:2:1:item:v}
  There is a dense open subset \(U\) of \(G\) such that, for every
  \(x\in U(k)\), the set of conjugates of \(H\) containing \(x\) is
  finite and nonempty; equivalently, \(\psi:X\to G\) is dominant and
  generically quasi-finite.
\item[(vi)] \SGARefTarget{sga3:xiii:theorem:2:1:item:vi}
  There is a dense open subset \(U\) of \(G\) such that every
  \(x\in U(k)\) is contained in a conjugate of \(H\); equivalently,
  \(\psi:X\to G\) is dominant.
\item[(vii)] \SGARefTarget{sga3:xiii:theorem:2:1:item:vii}
  There exists \(a\in H(k)\) such that the fixed subspace of
  \(\mathfrak g/\mathfrak h\) under
  \(\operatorname{ad}_{\mathfrak g/\mathfrak h}(a)\) is zero.
\end{enumerate}

\textit{Moreover, these conditions imply that \(H\) is its own
connected normalizer: \(N\) is smooth and
\(\dim H=\dim N\). They also imply that \(\psi:X\to G\) is
generically étale.}

\noindent\textit{Proof.}
By \SGARefLink{sga3:xiii:lemma:2:0}{2.0}, \(\dim X\leq\dim G\), with equality if and only if
\(\dim H=\dim N\), or equivalently if and only if \(H\) has finite
index in \(N\). The inequality \(\dim X\leq\dim G\) shows that
\(\psi\) is dominant if and only if it is dominant and generically
quasi-finite, and also if and only if it is generically quasi-finite
and \(\dim X=\dim G\). Since the last equality is equivalent, by
\SGARefLink{sga3:xiii:lemma:2:0}{2.0}, to \(H\) having finite index in \(N\), conditions \SGARefLink{sga3:xiii:theorem:2:1:item:vi}{(vi)}, \SGARefLink{sga3:xiii:theorem:2:1:item:v}{(v)}, and
\SGARefLink{sga3:xiii:theorem:2:1:item:iv-bis}{(iv bis)} are equivalent. The equivalence of \SGARefLink{sga3:xiii:theorem:2:1:item:iv}{(iv)} and \SGARefLink{sga3:xiii:theorem:2:1:item:iv-bis}{(iv bis)} is
immediate.

The equivalence of \SGARefLink{sga3:xiii:theorem:2:1:item:i}{(i)} and \SGARefLink{sga3:xiii:theorem:2:1:item:i-bis}{(i bis)} follows directly from the
definitions and is left to the reader. If \(H\) contains a Cartan
subgroup \(C\) of \(G\), it contains the maximal torus \(T\) of
\(C\), which is a maximal torus of \(G\). Since \(C\) is the
centralizer of \(T\), its Lie algebra \(\mathfrak c\) is
\[
  \mathfrak c=\mathfrak g^T
\]
(\SGARefLink{sga3:ii:theorem:5-2-3:item:ii}{II~5.2.3(ii)}). Hence
\(\mathfrak g^T=\mathfrak c\subset\mathfrak h\), and \SGARefLink{sga3:i:proposition:4-7-3}{I~4.7.3} gives
the relation\footnote{The French source has two evident typographical
slips in this sentence, writing \(\mathfrak n\) for \(\mathfrak c\)
and reversing the inclusion involving \(\mathfrak h\) and
\(\mathfrak g^T\). The mathematically required relation is restored
here.}
\SGARefTarget{sga3:xiii:theorem:2:1:equation:x}
\[
\tag{x}
  (\mathfrak g/\mathfrak h)^T=0.
\]
Conversely, suppose that \(H\) contains the maximal torus \(T\) and
that \SGARefLink{sga3:xiii:theorem:2:1:equation:x}{(x)} holds, so that \(\mathfrak h\supset\mathfrak c\). We claim
that \(H\) contains the centralizer \(C\) of \(T\), which will prove
the equivalence of \SGARefLink{sga3:xiii:theorem:2:1:item:i}{(i)} and \SGARefLink{sga3:xiii:theorem:2:1:item:ii}{(ii)}. This follows from \SGARefLink{sga3:xiii:lemma:2:1:1}{the next lemma}.

\medskip
\noindent\SGARefTarget{sga3:xiii:lemma:2:1:1}\textbf{Lemma 2.1.1.} ---
\textit{Let \(G\) be a smooth algebraic group over a field \(k\), let
\(T\) be a subgroup of multiplicative type of \(G\), let \(C\) be its
connected centralizer (equal to the centralizer of \(T\) when \(G\)
is connected and \(T\) is a torus, by \SGARefLink{sga3:xii:6.6:item:b}{XII~6.6(b)}), and let \(H\) be
a smooth subgroup of \(G\) containing \(T\). Then \(H\) contains
\(C\) if and only if its Lie algebra \(\mathfrak h\) contains the Lie
algebra \(\mathfrak c\) of \(C\).}

Indeed, \(\operatorname{Centr}_G(T)\) is smooth over \(k\) by \SGARefLink{sga3:xi:corollary:2:4}{XI~2.4},
and hence \(C\) is smooth over \(k\). Likewise,
\(\operatorname{Centr}_H(T)\) is smooth, and
\[
  \operatorname{Centr}_H(T)=\operatorname{Centr}_G(T)\cap H
\]
has Lie algebra \(\mathfrak h\cap\mathfrak c\). The hypothesis
therefore implies that this smooth subgroup has the same Lie algebra
as the smooth group \(\operatorname{Centr}_G(T)\), so it contains the
connected component \(C\) of the latter. Thus \(H\) contains \(C\).
\hfill\(\square\)

We now prove the equivalence of \SGARefLink{sga3:xiii:theorem:2:1:item:i}{(i)} and \SGARefLink{sga3:xiii:theorem:2:1:item:iii}{(iii)}. If \(H\) contains the
maximal torus \(T\) of \(G\), we must show that \(H\supset C\)
(equivalently, that \SGARefLink{sga3:xiii:theorem:2:1:equation:x}{(x)} holds) if and only if \(H\) has finite index
in its normalizer and the set of conjugates of \(H\) containing \(T\)
is finite. If \(H\supset C\), then \SGARefLink{sga3:xiii:theorem:2:1:equation:x}{(x)} implies, a fortiori,
\SGARefTarget{sga3:xiii:theorem:2:1:equation:xx}
\[
\tag{xx}
  (\mathfrak g/\mathfrak h)^H=0.
\]
The Lie algebra \(\mathfrak n\) is the inverse image of the left-hand
side under the canonical homomorphism
\(\mathfrak g\to\mathfrak g/\mathfrak h\) (II~5.2.3(i)).
Consequently, by \SGARefLink{sga3:xiii:lemma:2:0}{2.0}, \SGARefLink{sga3:xiii:theorem:2:1:equation:xx}{(xx)} says that \(H=N^0\); in particular, \(H\)
has finite index in its normalizer. Consider \SGARefLink{sga3:xiii:theorem:2:1:diagram:01}{the diagram of
subgroups}
\SGARefTarget{sga3:xiii:theorem:2:1:diagram:01}
\[
\begin{tikzcd}[column sep=large,row sep=large]
T \arrow[r] &
N(T)\cap H \arrow[r] \arrow[d] &
H \arrow[d] \\
{} &
N(T)\cap N(H) \arrow[r] \arrow[d] &
N(H) \\
{} &
N(T). &
\end{tikzcd}
\]
By the conjugacy theorem for maximal tori in \(H\) (\SGARefLink{sga3:xii:6.6:item:a}{XII~6.6(a)}),
every conjugate of \(H\) containing \(T\) is obtained from \(H\) by
conjugation by an element of \(N(T)(k)\). Thus the set of such
conjugates is in bijection with
\[
  \bigl(N(T)/(N(T)\cap N(H))\bigr)(k).
\]
Since \(H\supset C\), we have \(N(T)\cap H\supset C\), so the
preceding set is a quotient of \((N(T)/C)(k)\), which is finite.
This proves \SGARefLink{sga3:xiii:theorem:2:1:item:i}{(i)} \(\Rightarrow\) \SGARefLink{sga3:xiii:theorem:2:1:item:iii}{(iii)}.

Conversely, suppose \SGARefLink{sga3:xiii:theorem:2:1:item:iii}{(iii)}, so that
\(N(T)/(N(T)\cap N(H))\) and \(N(H)/H\) are finite. The conjugacy
theorem in \(H\) shows that the homomorphism
\[
  (N(T)\cap N(H))/(N(T)\cap H)\longrightarrow N(H)/H
\]
induced by \SGARefLink{sga3:xiii:theorem:2:1:diagram:01}{the preceding diagram} is bijective on \(k\)-valued points
(indeed, it is an isomorphism). Since the latter quotient is finite,
so is the former. Hence \(N(T)\cap H\) has finite index in \(N(T)\)
and therefore contains \(C=N(T)^0\). Thus \(H\supset C\), and \SGARefLink{sga3:xiii:theorem:2:1:item:i}{(i)},
\SGARefLink{sga3:xiii:theorem:2:1:item:i-bis}{(i bis)}, \SGARefLink{sga3:xiii:theorem:2:1:item:ii}{(ii)}, and \SGARefLink{sga3:xiii:theorem:2:1:item:iii}{(iii)} are equivalent.

We next prove \SGARefLink{sga3:xiii:theorem:2:1:item:ii}{(ii)} \(\Rightarrow\) \SGARefLink{sga3:xiii:theorem:2:1:item:vii}{(vii)}. Both conditions are
invariant under an extension \(k\to k'\) with \(k'\) algebraically
closed. We may therefore assume that \(k\) has infinite
transcendence degree over its prime subfield. There is then an
element \(a\in T(k)\) such that the subgroup of \(T(k)\) generated by
\(a\) is Zariski dense in \(T\). It follows that
\[
  (\mathfrak g/\mathfrak h)^T
  =(\mathfrak g/\mathfrak h)^{\operatorname{ad}(a)}.
\]
The left-hand side is zero by hypothesis, which proves \SGARefLink{sga3:xiii:theorem:2:1:item:vii}{(vii)}.

The implication \SGARefLink{sga3:xiii:theorem:2:1:item:vii}{(vii)} \(\Rightarrow\) \SGARefLink{sga3:xiii:theorem:2:1:item:vi}{(vi)} is contained in \SGARefLink{sga3:xiii:corollary:2:2}{the
following result}, which sharpens \SGARefLink{sga3:xiii:theorem:2:1}{2.1}.

\medskip
\noindent\SGARefTarget{sga3:xiii:corollary:2:2}\textbf{Corollary 2.2.} ---
\textit{Let \(G\) be a smooth algebraic group over a field \(k\), let
\(H\) be a smooth algebraic subgroup, let \(N\) be its normalizer in
\(G\), and let}
\[
  \varphi:G\times H\longrightarrow G,\qquad
  \varphi(g,h)=\operatorname{ad}(g)h=ghg^{-1}.
\]
\textit{Let \(\psi:X=G\times^N H\to G\) be the morphism obtained from
\(\varphi\) by passage to the quotient (cf. \SGARefLink{sga3:xiii:section:1}{Section 1}), and let
\(a\in H(k)\). The following conditions are equivalent:}

\begin{enumerate}
\item[(i)] \SGARefTarget{sga3:xiii:corollary:2:2:item:i}
  \(\varphi\) is smooth at \((e,a)\).
\item[(ii)] \SGARefTarget{sga3:xiii:corollary:2:2:item:ii}
  \(\psi\) is étale at \((\bar e,a)\), and \(N\) is smooth over \(k\).
\item[(iii)] \SGARefTarget{sga3:xiii:corollary:2:2:item:iii}
  \((\mathfrak g/\mathfrak h)^{\operatorname{ad}(a)}=0\), where
  \(\mathfrak g\) and \(\mathfrak h\) are the Lie algebras of \(G\)
  and \(H\).
\end{enumerate}

\textit{These conditions imply \(H^0=N^0\).}

The smoothness at a \(k\)-rational point of \(\varphi\), a morphism
between smooth \(k\)-preschemes, is equivalent to the surjectivity of
the tangent map at that point. Under the usual identifications of
the tangent spaces of \(G\) and \(H\) with their Lie algebras, a
direct calculation gives
\SGARefTarget{sga3:xiii:corollary:2:2:display:tangent-map}
\[
  d\varphi(\xi,\eta)
  =\bigl(\operatorname{id}-\operatorname{ad}(a)\bigr)\xi+\eta,
\]
as a map from \(\mathfrak g\times\mathfrak h\) to \(\mathfrak g\).
Its surjectivity is therefore equivalent to the surjectivity of
\(\operatorname{id}-\operatorname{ad}(a)\) on
\(\mathfrak g/\mathfrak h\), which is precisely \SGARefLink{sga3:xiii:corollary:2:2:item:iii}{(iii)}.

Condition \SGARefLink{sga3:xiii:corollary:2:2:item:iii}{(iii)} implies, a fortiori,
\((\mathfrak g/\mathfrak h)^H=0\), and hence \(H^0=N^0\) by \SGARefLink{sga3:xiii:lemma:2:0}{2.0}
(the connectedness hypothesis made at the beginning of the section
is not needed there). It follows that \(N\) is smooth and that
\(\dim H=\dim N\), whence \(\dim X=\dim G\). Since
\(q:G\times H\to X\) is flat and
\(\psi\circ q=\varphi\)\footnote{The source reverses this composition,
printing \(\psi=\varphi\circ q\). The displayed morphisms and the
definition of \(\varphi\) require \(\varphi=\psi\circ q\).} is smooth
at \((e,a)\), the morphism \(\psi\)
is smooth at \(q(e,a)=(\bar e,a)\), and hence is étale there by
dimension. We have proved
\(\SGARefLink{sga3:xiii:corollary:2:2:item:i}{\mathrm{(i)}}\Longleftrightarrow\SGARefLink{sga3:xiii:corollary:2:2:item:iii}{\mathrm{(iii)}}
\Longrightarrow\SGARefLink{sga3:xiii:corollary:2:2:item:ii}{\mathrm{(ii)}}\). Conversely, \SGARefLink{sga3:xiii:corollary:2:2:item:ii}{(ii)} implies \SGARefLink{sga3:xiii:corollary:2:2:item:i}{(i)},
because the smoothness of \(N\) implies that of \(q\).
\hfill\(\square\)

In particular, under condition \SGARefLink{sga3:xiii:theorem:2:1:item:vii}{(vii)} of \SGARefLink{sga3:xiii:theorem:2:1}{2.1} the morphism
\(\varphi\) is smooth at \((e,a)\), so its image contains a nonempty
open subset of \(G\). Thus \(\psi\) is dominant, proving
\SGARefLink{sga3:xiii:theorem:2:1:item:vii}{(vii)} \(\Rightarrow\) \SGARefLink{sga3:xiii:theorem:2:1:item:vi}{(vi)}.

It remains to prove \SGARefLink{sga3:xiii:theorem:2:1:item:vi}{(vi)} \(\Rightarrow\) \SGARefLink{sga3:xiii:theorem:2:1:item:i}{(i)}. Suppose first that
\(G\) is affine. Let \(U\) be a nonempty open subset of \(G\) such
that every \(x\in U(k)\) is contained in a conjugate of \(H\), and
let \(C\) be a Cartan subgroup of \(G\). Applying the already proved
implication \SGARefLink{sga3:xiii:theorem:2:1:item:i}{(i)} \(\Rightarrow\) \SGARefLink{sga3:xiii:theorem:2:1:item:iv}{(iv)} with \(C\) in place of \(H\)
(Borel's ``density theorem''), we find a conjugate of \(C\) meeting
\(U\). We may therefore assume that \(U\cap C\neq\varnothing\);
equivalently, there is a nonempty open subset \(V\) of \(C\) such
that every \(x\in V(k)\) belongs to a conjugate of \(H\).

Write
\SGARefTarget{sga3:xiii:theorem:2:1:display:cartan-decomposition}
\[
  C=T\cdot C_u,
\]
where \(T\) is the maximal torus of \(C\), hence a maximal torus of
\(G\), and \(C_u\) is the unipotent part of \(C\); \(T\) lies in the
center of \(C\) (\emph{BIBLE} 6, Theorem 2). We may again assume that
\(k\) has infinite transcendence degree over its prime subfield.
Choose \(t\in T(k)\) in the projection of \(V\) onto \(T\), a
nonempty open subset of \(T\), so that
\(tC_u\cap V\neq\varnothing\), and choose it so that it generates a
Zariski-dense subgroup of \(T\). Every algebraic subgroup of \(G\)
containing a product \(tu\), with \(t\in T(k)\) and
\(u\in C_u(k)\), contains both factors (\emph{BIBLE} 4,
Theorem 3). Taking \(tu\in V(k)\), we conclude that a conjugate of
\(H\) contains \(t\), and therefore contains \(T\). After replacing
\(H\) by that conjugate, we may assume
\[
  T\subset H.
\]

Let \(W\) be the inverse image in \(C_u\) of \(V\) under
\(u\mapsto tu\). For every \(x\in(T\cdot W)(k)\), there is then a conjugate
of \(H\) containing both \(T\) and \(x\). As observed above, every
such conjugate is of the form \(\operatorname{int}(g)H\) for some
\(g\in N(T)(k)\). Consider
\SGARefTarget{sga3:xiii:theorem:2:1:display:f}
\[
  f:N(T)\times H\longrightarrow G,\qquad
  f(g,h)=\operatorname{int}(g)h=ghg^{-1}.
\]
The image of \(f\) contains \(TW\). Since \(C\) has finite index in
\(N(T)\), the latter is a finite union of translates \(Cg_i\), with
\(g_i\in N(T)(k)\). It follows that some dense open subset \(V'\) of
\(C=T C_u\) is contained in the image under \(f\) of
\((Cg_i)\times H\). Replacing \(H\) by
\(\operatorname{int}(g_i)H\), we may assume \(g_i=e\), so that
\[
  f(C\times H)\supset V'.
\]
Thus for every \(u\in V'(k)\) there are \(v\in C(k)\) and
\(h\in H(k)\) such that
\[
  v^{-1}hv=u,\qquad\text{hence}\qquad vuv^{-1}\in H(k).
\]
Set
\[
  C'=C\cap H=\operatorname{Centr}_H(T).
\]
Since \(vuv^{-1}\) also belongs to \(C\), it lies in \(C'\), and
hence \(u\in\operatorname{int}(v^{-1})C'\). Therefore the union of
the conjugates of \(C'\) in \(C\), by elements of \(C(k)\), is
dense. As in the preceding argument for the pair \((G,H)\), now
applied to \((C,C')\), this implies that \(C'\) has finite index in
its normalizer in \(C\). By the lemma of \emph{BIBLE} 7, \SGARefLink{sga3:xiii:section:1}{Section 1},
we obtain \(C'=C\), and consequently \(C\subset H\).\footnote{The
source prints \(C=H\); only the inclusion \(C\subset H\) follows and
is used.} This proves
\SGARefLink{sga3:xiii:theorem:2:1:item:vi}{(vi)} \(\Rightarrow\) \SGARefLink{sga3:xiii:theorem:2:1:item:i}{(i)} when \(G\) is affine.

In the general case we argue by induction on \(n=\dim G\), the
assertion being trivial for \(n=0\). Let \(Z\) be the center of
\(G\). There are two cases.

\begin{enumerate}
\item[\(1^\circ\)] \SGARefTarget{sga3:xiii:theorem:2:1:proof:case:1} If \(\dim(Z\cap H)>0\), set
  \(G'=G/(Z\cap H)\). Then \(\dim G'<n\), and condition \SGARefLink{sga3:xiii:theorem:2:1:item:vi}{(vi)} for
  \(H\) implies \SGARefLink{sga3:xiii:theorem:2:1:item:vi}{the same condition} for its image \(H'\) in \(G'\).
  By induction, \(H'\) contains a Cartan subgroup \(C'\) of \(G'\).
  Hence \(H\) contains the inverse image \(C\) of \(C'\), which is a
  Cartan subgroup by \SGARefLink{sga3:xii:theorem:6-6:item:e}{XII~6.6(e)}.

\item[\(2^\circ\)] \SGARefTarget{sga3:xiii:theorem:2:1:proof:case:2} If \(\dim(Z\cap H)=0\), the canonical morphism
  \(H\to G/Z\) is finite. Since \(G/Z\) is affine by \SGARefLink{sga3:xii:6.1}{XII~6.1}, \(H\)
  is affine. Every homomorphism from \(H\) to an abelian variety is
  therefore zero (indeed, every morphism of preschemes from \(H\) to
  an abelian variety is constant). This follows from the fact that a
  smooth connected affine algebraic group over an algebraically
  closed field is a rational variety, or simply from the fact that
  it is the union of its Borel subgroups (\emph{BIBLE} 6,
  Theorem 5(b)); the structure theorems \emph{BIBLE} 6.2 and 6.3
  show readily that a smooth connected solvable affine group is a
  rational variety.

  Apply Chevalley's structure theorem to \(G\): the group \(G\) is
  an extension of an abelian variety \(A\) by a smooth affine group.
  Since \(H\) is affine, its image in \(A\) is zero. On the other
  hand, that image must be all of \(A\), because the union of its
  conjugates in \(A\) is dense. Hence \(A=0\), so \(G\) is affine,
  reducing us to the case already proved.
\end{enumerate}

This completes the proof of \SGARefLink{sga3:xiii:theorem:2:1}{2.1}.

\medskip
\noindent\SGARefTarget{sga3:xiii:corollary:2:3}\textbf{Corollary 2.3.} ---
\textit{Assume that the equivalent conditions of \SGARefLink{sga3:xiii:theorem:2:1}{2.1} hold.}

\begin{enumerate}
\item[\textit{a)}] \SGARefTarget{sga3:xiii:corollary:2:3:item:a}
  Let \(k(X)\) and \(k(G)\) be the fields of rational functions of
  \(X\) and \(G\), respectively. Then \(k(X)\) is a finite separable
  extension of \(k(G)\). Denote its degree by \(d\).

\item[\textit{b)}] \SGARefTarget{sga3:xiii:corollary:2:3:item:b}
  Let \(T\) be a maximal torus of \(G\) contained in \(H\), which
  exists by \SGARefLink{sga3:xiii:theorem:2:1:item:ii}{2.1(ii)}, and let \(C\) be the corresponding Cartan
  subgroup of \(G\). Then \(C\subset H\). Moreover,
  \(\operatorname{Norm}_G(T)\) is a smooth subgroup of \(G\), and
  \[
    \operatorname{Norm}_G(T)\cap\operatorname{Norm}_G(H)
    =
    \operatorname{Norm}_{\operatorname{Norm}_G(H)}(T)
  \]
  is a smooth subgroup of it, of finite index equal to the degree
  \(d\) defined in \SGARefLink{sga3:xiii:corollary:2:3:item:a}{(a)}. The number of conjugates of \(H\) containing
  a given maximal torus, or a given Cartan subgroup, is \(d\).

\item[\textit{c)}] \SGARefTarget{sga3:xiii:corollary:2:3:item:c}
  Let \(U\) be the largest open subset of \(G\) such that
  \(\psi:X\to G\) induces a finite étale morphism
  \[
    \psi^{-1}(U)\longrightarrow U.
  \]
  Then \(U\) is dense. For \(g\in G(k)\), the following are
  equivalent:
  \begin{enumerate}
  \item \SGARefTarget{sga3:xiii:corollary:2:3:item:c:condition:i} \(g\in U(k)\);
  \item \SGARefTarget{sga3:xiii:corollary:2:3:item:c:condition:ii} exactly \(d\) conjugates of \(H\) contain \(g\);
  \item \SGARefTarget{sga3:xiii:corollary:2:3:item:c:condition:iii} there are at least \(d\) distinct conjugates \(H_i\) of \(H\)
    containing \(g\), and for every \(i\),
    \[
      (\mathfrak g/\mathfrak h_i)^{\operatorname{ad}(g)}=0,
      \qquad \mathfrak h_i=\operatorname{Lie}(H_i).
    \]
  \end{enumerate}
\end{enumerate}

Assertion \SGARefLink{sga3:xiii:corollary:2:3:item:a}{(a)} follows from the generic étaleness of \(\psi\), stated
at the end of \SGARefLink{sga3:xiii:theorem:2:1}{2.1}. This also shows that the open set \(U\) in \SGARefLink{sga3:xiii:corollary:2:3:item:c}{(c)} is
nonempty, hence dense, and gives the two stated characterizations of
its \(k\)-points: here one uses that \(\psi\) is separated, that
\(X\) is integral, that \(G\) is integral and normal (SGA 1,
I~10.11), and that
\((\mathfrak g/\mathfrak h_i)^{\operatorname{ad}(g)}=0\) means that
\(\psi\) is étale at the point \(x_i\in\psi^{-1}(g)\) corresponding
to \(H_i\).

If \(H\) contains the maximal torus \(T\) of \(G\), the centralizers
of \(T\) in \(H\) and \(G\) have the same dimension and are smooth
and connected (\SGARefLink{sga3:xii:theorem:6-6:item:b}{XII~6.6(b)}); they are therefore equal, proving that
\(C\subset H\).\footnote{The source prints the tautological
\(C\subset G\); equality of the two centralizers gives
\(C\subset H\).} The normalizer of \(T\) in any smooth group
containing it is smooth (\SGARefLink{sga3:xi:corollary:2:4:bis}{XI~2.4 bis}), so
\(\operatorname{Norm}_G(T)\) and
\(\operatorname{Norm}_{\operatorname{Norm}_G(H)}(T)\) are smooth.
Recall also that \(N=\operatorname{Norm}_G(H)\) is smooth by the
last assertion of \SGARefLink{sga3:xiii:theorem:2:1}{2.1}. Furthermore,
\(\operatorname{Norm}_N(T)\) contains \(C\), which has finite index
in \(\operatorname{Norm}_G(T)\); hence
\(\operatorname{Norm}_N(T)\) itself has finite index in
\(\operatorname{Norm}_G(T)\).

The conjugacy theorem for maximal tori in \(H\) shows that this index
is the number of conjugates of \(H\) containing \(T\), equivalently
the number containing \(C\). The union of the conjugates of \(C\) in
\(G\) is dense, by \SGARefLink{sga3:xiii:theorem:2:1:item:i}{2.1(i)} \(\Rightarrow\) \SGARefLink{sga3:xiii:theorem:2:1:item:vi}{(vi)} applied to \(C\) in
place of \(H\), and \(U\) is stable under inner automorphisms. Hence
\(C\cap U\neq\varnothing\). Arguing as in the proof of
\SGARefLink{sga3:xiii:theorem:2:1:item:vi}{2.1(vi)} \(\Rightarrow\) \SGARefLink{sga3:xiii:theorem:2:1:item:ii}{(ii)}, and making a harmless extension of the
base field if necessary, one finds
\(g\in(C\cap U)(k)\) such that every conjugate of \(H\) containing
\(g\) contains \(T\), and hence \(C\). Thus the conjugates of \(H\)
containing \(C\) are precisely those containing \(g\); since
\(g\in U(k)\), their number is \(d\). This proves \SGARefLink{sga3:xiii:corollary:2:3:item:b}{(b)}.

In fact, the argument shows that the set of conjugates of \(H\)
containing \(T\) is a homogeneous set under the group of rational
points of
\[
  W_G(T)=
  \operatorname{Norm}_G(T)/\operatorname{Centr}_G(T).
\]
In particular,
\[
  d\leq\text{the order of the Weyl group of \(G\)}.
\]

\medskip
\noindent\SGARefTarget{sga3:xiii:corollary:2:4}\textbf{Corollary 2.4.} ---
\textit{With the notation of \SGARefLink{sga3:xiii:theorem:2:1}{2.1}, the following conditions are
equivalent:}

\begin{enumerate}
\item[(i)] \SGARefTarget{sga3:xiii:corollary:2:4:item:i}
  \(\psi:X\to G\) is birational.
\item[(ii)] \SGARefTarget{sga3:xiii:corollary:2:4:item:ii}
  Exactly one conjugate of \(H\) contains a given Cartan subgroup of
  \(G\).
\item[(iii)] \SGARefTarget{sga3:xiii:corollary:2:4:item:iii}
  \(H\) contains a Cartan subgroup \(C\) of \(G\), and
  \(\operatorname{Norm}_G(H)\supset\operatorname{Norm}_G(C)\).
\item[(iv)] \SGARefTarget{sga3:xiii:corollary:2:4:item:iv}
  There is a nonempty open subset \(V\) of \(G\) such that every
  \(g\in V(k)\) belongs to exactly one conjugate of \(H\).
\end{enumerate}

This follows immediately from \SGARefLink{sga3:xiii:theorem:2:1}{2.1} and \SGARefLink{sga3:xiii:corollary:2:3}{2.3}.

\medskip
\noindent\SGARefTarget{sga3:xiii:corollary:2:5}\textbf{Corollary 2.5.} ---
\textit{Assume that the conditions of \SGARefLink{sga3:xiii:corollary:2:4}{2.4} hold, and let
\(g\in G(k)\). The following conditions are equivalent:}

\begin{enumerate}
\item[(i)] \SGARefTarget{sga3:xiii:corollary:2:5:item:i}
  \(g\in U(k)\), where \(U\) is defined in \SGARefLink{sga3:xiii:corollary:2:3:item:c}{2.3(c)}; that is, \(g\) is
  contained in exactly one conjugate of \(H\).
\item[(ii)] \SGARefTarget{sga3:xiii:corollary:2:5:item:ii}
  The set of conjugates of \(H\) containing \(g\) is finite and
  nonempty.
\item[(iii)] \SGARefTarget{sga3:xiii:corollary:2:5:item:iii}
  The scheme \(\psi^{-1}(g)\), the ``scheme of conjugates of \(H\)
  containing \(g\),'' has an isolated point.
\item[(iv)] \SGARefTarget{sga3:xiii:corollary:2:5:item:iv}
  There is a conjugate \(H'\) of \(H\) containing \(g\) such that
  \[
    (\mathfrak g/\mathfrak h')^{\operatorname{ad}(g)}=0,
    \qquad \mathfrak h'=\operatorname{Lie}(H').
  \]
\end{enumerate}

\textit{Finally, \(U\) is also the largest open subset of \(G\) on
which \(\psi\) induces an isomorphism}
\[
  \psi^{-1}(U)\xrightarrow{\sim}U.
\]

The equivalence of \SGARefLink{sga3:xiii:corollary:2:5:item:i}{(i)}, \SGARefLink{sga3:xiii:corollary:2:5:item:ii}{(ii)}, and \SGARefLink{sga3:xiii:corollary:2:5:item:iii}{(iii)}, together with the last
assertion, follows from the ``Main Theorem'' applied to the
birational morphism \(\psi:X\to G\), since \(G\) is normal. Their
equivalence with \SGARefLink{sga3:xiii:corollary:2:5:item:iv}{(iv)} follows immediately from \SGARefLink{sga3:xiii:corollary:2:2}{Corollary~2.2}, which
characterizes the points of \(X\) at which \(\psi\) is
étale.\footnote{The source cites the final assertion of \SGARefLink{sga3:xiii:theorem:2:1}{2.1} here;
the pointwise étaleness criterion is \SGARefLink{sga3:xiii:corollary:2:2}{Corollary 2.2}.}

\medskip
\noindent\SGARefTarget{sga3:xiii:theorem:2:6}\textbf{Theorem 2.6.} ---
\textit{Let \(G\) be a smooth connected algebraic group over an
algebraically closed field \(k\), let \(C\) be a Cartan subgroup
associated with a maximal torus \(T\), and put}
\[
  N=\operatorname{Norm}_G(C)=\operatorname{Norm}_G(T)
  \qquad\text{(cf. XII~8.4).}
\]
\textit{Let \(X=G\times^N C\), where \(N\) acts on the left factor
\(G\) by right translations and on the right factor \(C\) by inner
automorphisms, and let \(\psi:X\to G\) be the canonical morphism.}

\begin{enumerate}
\item[\textit{a)}] \SGARefTarget{sga3:xiii:theorem:2:6:item:a}
  \(\psi\) is birational.

\item[\textit{b)}] \SGARefTarget{sga3:xiii:theorem:2:6:item:b}
  Let \(U\) be the largest open subset of \(G\) such that \(\psi\)
  induces an isomorphism
  \[
    \psi^{-1}(U)\xrightarrow{\sim}U
  \]
  (cf. \SGARefLink{sga3:xiii:corollary:2:5}{2.5}), and let
  \SGARefTarget{sga3:xiii:theorem:2:6:display:rho}
  \[
    \rho=\rho_\nu(G)=\dim C
  \]
  be the nilpotent rank of \(G\). For every \(g\in G(k)\), the
  multiplicity of the eigenvalue \(1\) in the action of
  \(\operatorname{ad}(g)\) on \(\mathfrak g\) is at least \(\rho\);
  it equals \(\rho\) if and only if \(g\in U(k)\).
\end{enumerate}

\noindent\textit{Proof.}
\SGARefLink{sga3:xiii:theorem:2:1:item:i}{Condition 2.1(i)} holds, so \SGARefLink{sga3:xiii:corollary:2:4:item:iii}{2.4(iii)} \(\Rightarrow\) \SGARefLink{sga3:xiii:corollary:2:4:item:i}{(i)} proves \SGARefLink{sga3:xiii:theorem:2:6:item:a}{(a)}.
In \emph{BIBLE} 7, when \(G\) is affine, the points of \(U(k)\) are
called the \emph{regular points} of \(G(k)\). We shall use the same
terminology and call \(U\) the \emph{open subset of regular points}
of \(G\). Notice that the proof in \emph{BIBLE} that this set is
itself open is incorrect; \SGARefLink{sga3:xiii:expose}{the present Exposé} establishes its
openness under more general hypotheses.

To prove \SGARefLink{sga3:xiii:theorem:2:6:item:b}{(b)}, for each \(g\in G(k)\) introduce \SGARefLink{sga3:xiii:theorem:2:6:display:characteristic-polynomial}{the characteristic
polynomial}
\SGARefTarget{sga3:xiii:theorem:2:6:display:characteristic-polynomial}
\[
  P(\operatorname{ad}(g),t)
  =t^n+c_1(g)t^{n-1}+\cdots+c_n(g).
\]
Replacing \(k\) by an arbitrary \(k\)-algebra shows immediately that
the functions \(c_i(g)\) arise from uniquely determined sections
\[
  c_i\in\Gamma(G,\mathcal O_G).
\]

Suppose that \(g\in G(k)\) belongs to a Cartan subgroup, which we may
take to be \(C\); this includes the case of a regular element. By
\SGARefLink{sga3:xiii:corollary:2:5:item:iv}{2.5(iv)},
\[
  (\mathfrak g/\mathfrak c)^{\operatorname{ad}(g)}=0
\]
if and only if \(g\) is regular, where
\(\mathfrak c=\operatorname{Lie}(C)\). Since \(C\) is nilpotent,
\(\operatorname{ad}_{\mathfrak c}(g)\) has only the eigenvalue \(1\).
It follows that the multiplicity of the eigenvalue \(1\) in
\(\operatorname{ad}_{\mathfrak g}(g)\) is at least \(\rho\), and is
exactly \(\dim C=\rho\) if and only if \(g\) is regular.

In particular, \SGARefLink{sga3:xiii:theorem:2:6:display:characteristic-polynomial}{the characteristic polynomial above} is divisible by
\((t-1)^\rho\). Divisibility by \((t-1)^\rho\) is expressed by
linear relations, with integer coefficients, among the coefficients
of the polynomial. These relations hold for \(g\in U(k)\), where
\(U\) is dense; since \(G\) is reduced, they therefore hold
everywhere. Thus\footnote{The French source prints \(c_0\) in place
of the omitted tail \(+\cdots+c_n\) on the left-hand side of
\SGARefLink{sga3:xiii:theorem:2:6:equation:dagger}{\((\dagger)\)}. \SGARefLink{sga3:xiii:theorem:2:6:display:characteristic-polynomial}{The characteristic polynomial displayed immediately
above} determines the correction made here.}
\SGARefTarget{sga3:xiii:theorem:2:6:equation:dagger}
\[
\tag{\(\dagger\)}
  t^n+c_1t^{n-1}+\cdots+c_n
  =(t-1)^\rho
    \bigl(t^{n-\rho}+b_1t^{n-\rho-1}+\cdots+b_{n-\rho}\bigr)
\]
in the polynomial ring over \(\Gamma(G,\mathcal O_G)\). In
particular, for every field extension \(K/k\) and every
\(g\in G(K)\), the eigenvalue \(1\) of \(\operatorname{ad}(g)\) has
multiplicity at least \(\rho\).

We have already seen that equality holds when \(g\) is regular; let
us prove the converse. Suppose first that \(G\) is affine, and write
\[
  g=g_sg_u
\]
as the product of its semisimple and unipotent parts
(\emph{BIBLE} 4, Section 4). Then
\[
  \operatorname{ad}(g)
  =\operatorname{ad}(g_s)\operatorname{ad}(g_u)
\]
is the corresponding decomposition of \(\operatorname{ad}(g)\)
(loc. cit., corollary to Theorem 3). Hence
\(\operatorname{ad}(g)\) and \(\operatorname{ad}(g_s)\) have the same
eigenvalues, counted with multiplicity; in particular, the
eigenvalue \(1\) has the same multiplicity in both.

By \emph{BIBLE} 7, Theorem 2, Corollary 1, the element \(g\) is
regular if and only if \(g_s\) is regular. We may therefore assume
that \(g=g_s\) is semisimple. It is then contained in a maximal torus
by \emph{BIBLE} 6, Theorem 5(c), and hence in a Cartan subgroup, a
case already treated. This proves \SGARefLink{sga3:xiii:theorem:2:6:item:b}{(b)} when \(G\) is affine.

In general, let \(Z=\operatorname{Centr}(G)_{\mathrm{red}}\). By
\SGARefLink{sga3:xii:6.6:item:e}{XII~6.6(e)}, the Cartan subgroups of \(G\) are the inverse images of
those of \(G'=G/Z\). Thus \(g\) is regular in \(G\) if and only if
its image \(g'\) is regular in \(G'\). Since \(Z\) is smooth, the Lie
algebra of \(G'\) is
\(\mathfrak g'=\mathfrak g/\mathfrak z\), where
\(\mathfrak z=\operatorname{Lie}(Z)\), and
\(\operatorname{ad}(g')\) is the endomorphism induced by
\(\operatorname{ad}(g)\) on \(\mathfrak g/\mathfrak z\).
Consequently, the multiplicity of the eigenvalue \(1\) in
\(\operatorname{ad}(g)\) is
\[
  d=\dim Z
\]
plus its multiplicity in \(\operatorname{ad}(g')\). The former
therefore equals the nilpotent rank of \(G\) if and only if the latter
equals the nilpotent rank of \(G'\). We are reduced to \(G'\), which
is affine by \SGARefLink{sga3:xii:6.1}{XII~6.1}, and that case has already been proved. This
completes the proof of \SGARefLink{sga3:xiii:theorem:2:6}{2.6}.

\medskip
\noindent\SGARefTarget{sga3:xiii:corollary:2:7}\textbf{Corollary 2.7.} ---
\textit{With the notation of the preceding proof, let}
\[
  b=1+b_1+\cdots+b_{n-\rho}\in\Gamma(G,\mathcal O_G).
\]
\textit{Then the open subset of regular points of \(G\) is}
\[
  U=G_b,
\]
\textit{the set of points of \(G\) at which \(b\) is invertible. In
particular, \(U\) is an affine open subset when \(G\) is affine.}

\medskip
\noindent\SGARefTarget{sga3:xiii:corollary:2:8}\textbf{Corollary 2.8.} ---
\textit{Let \(H\) be a smooth connected algebraic subgroup of \(G\)
containing a Cartan subgroup of \(G\).}

\begin{enumerate}
\item[\textit{a)}] \SGARefTarget{sga3:xiii:corollary:2:8:item:a}
  Let \(C\) be an algebraic subgroup of \(H\). Then \(C\) is a
  Cartan subgroup of \(H\) if and only if it is a Cartan subgroup of
  \(G\).

\item[\textit{b)}] \SGARefTarget{sga3:xiii:corollary:2:8:item:b}
  Let \(g\in G(k)\), and let \(d\) be the integer introduced in \SGARefLink{sga3:xiii:corollary:2:3}{2.3}.
  Then \(g\) is a regular point of \(G\) if and only if there are
  exactly \(d\) conjugates \(H_i\) of \(H\) containing \(g\), and
  \(g\) is a regular element of every \(H_i\). Equivalently, there
  are at most \(d\) conjugates of \(H\) containing \(g\), and \(g\)
  is regular in one of them. When these conditions hold, if \(C\) is
  the unique Cartan subgroup of \(G\) containing \(g\), then the
  conjugates of \(H\) containing \(g\) are precisely the conjugates
  of \(H\) containing \(C\).

\item[\textit{c)}] \SGARefTarget{sga3:xiii:corollary:2:8:item:c}
  Let \(g\in H(k)\). Then \(g\) is regular in \(G\) if and only if it
  is regular in \(H\) and
  \[
    (\mathfrak g/\mathfrak h)^{\operatorname{ad}(g)}=0.
  \]
\end{enumerate}

\noindent\textit{Proof.}
\SGARefTarget{sga3:xiii:corollary:2:8:proof:item:a}
\textit{a)} Under either hypothesis on \(C\), its unique maximal
torus \(T\) is a maximal torus of both \(H\) and \(G\), since \(H\)
has the same reductive rank as \(G\). The smooth connected groups
\(\operatorname{Centr}_H(T)\subset\operatorname{Centr}_G(T)\) have
the same dimension, and therefore are equal. Hence \(C\) is equal to
one of these two centralizers if and only if it is equal to the
other, proving \SGARefLink{sga3:xiii:corollary:2:8:item:a}{(a)}.

\SGARefTarget{sga3:xiii:corollary:2:8:proof:item:b}
\textit{b)} Suppose first that \(g\) is regular in \(G\), and let
\(C\) be the unique Cartan subgroup of \(G\) containing \(g\). By
\SGARefLink{sga3:xiii:corollary:2:3:item:b}{2.3(b)}, exactly \(d\) conjugates \(H_i\) of \(H\) contain \(C\).
Since
\[
  (\mathfrak g/\mathfrak c)^{\operatorname{ad}(g)}=0,
\]
the endomorphism \(\operatorname{ad}(g)\) has no eigenvalue \(1\) on
\(\mathfrak g/\mathfrak c\), and a fortiori
\[
  (\mathfrak g/\mathfrak h_i)^{\operatorname{ad}(g)}=0.
\]
By \SGARefLink{sga3:xiii:corollary:2:3:item:c}{2.3(c)}, exactly \(d\) conjugates of \(H\) contain \(g\), namely
the \(H_i\). A Cartan subgroup of such an \(H_i\) containing \(g\) is
a Cartan subgroup of \(G\) by \SGARefLink{sga3:xiii:corollary:2:8:item:a}{(a)}, and hence is \(C\). Thus \(g\) is
regular in every \(H_i\).

Conversely, suppose that at most \(d\) conjugates \(H_i\) of \(H\)
contain \(g\), and that \(g\) is regular in one of them, which we may
take to be \(H\). Since \(g\) is regular in \(H\), it belongs to a
unique Cartan subgroup \(C\) of \(H\); by \SGARefLink{sga3:xiii:corollary:2:8:item:a}{(a)}, \(C\) is a Cartan
subgroup of \(G\). Let \(C'\) be any Cartan subgroup of \(G\)
containing \(g\). By \SGARefLink{sga3:xiii:corollary:2:3:item:b}{2.3(b)}, exactly \(d\) conjugates of \(H\)
contain \(C'\). They all contain \(g\), so they are precisely the
\(H_i\); in particular, \(H\) contains \(C'\). Thus \(C\) and \(C'\)
are two Cartan subgroups of \(H\) containing the same regular element
\(g\) of \(H\), and therefore \(C=C'\). Hence \(g\) is regular in
\(G\).\footnote{The source prints ``regular in \(C\)'' at this
conclusion; the argument proves regularity in \(G\).}
\hfill\(\square\)

\SGARefTarget{sga3:xiii:corollary:2:8:proof:item:c}
\textit{c)} For an endomorphism \(u\) of a finite-dimensional vector
space, let \(\nu(u)\) denote the nullity of \(\operatorname{id}-u\). Then
\SGARefTarget{sga3:xiii:corollary:2:8:display:nullity}
\[
  \nu\bigl(\operatorname{ad}(g)_{\mathfrak g}\bigr)
  =
  \nu\bigl(\operatorname{ad}(g)_{\mathfrak h}\bigr)
  +
  \nu\bigl(\operatorname{ad}(g)_{\mathfrak g/\mathfrak h}\bigr).
\]
The two terms on the right are respectively at least the nilpotent
rank of \(H\), which equals the nilpotent rank \(\rho\) of \(G\) by
\SGARefLink{sga3:xiii:corollary:2:8:item:a}{(a)}, and at least \(0\). Hence
\(\nu(\operatorname{ad}(g)_{\mathfrak g})=\rho\) if and only if
\[
  \nu\bigl(\operatorname{ad}(g)_{\mathfrak h}\bigr)=\rho,
  \qquad
  \nu\bigl(\operatorname{ad}(g)_{\mathfrak g/\mathfrak h}\bigr)=0.
\]
This says precisely that \(g\) is regular in \(G\) if and only if it
is regular in \(H\) and
\(\operatorname{ad}(g)_{\mathfrak g/\mathfrak h}\) has no nonzero
invariants. \hfill\(\square\)

\medskip
\noindent\SGARefTarget{sga3:xiii:remark:2:9}\textbf{Remarks 2.9.} ---
In the statement of \SGARefLink{sga3:xiii:theorem:2:1}{2.1}, condition \SGARefLink{sga3:xiii:theorem:2:1:item:iii}{(iii)} cannot be weakened by
assuming only that \(H\) contains a maximal torus and has finite
index in its normalizer, even if one also requires that normalizer to
be smooth, that is, that \(H=N^0\), and even when \(G\) is affine and
solvable. An example is the group \(G\) of matrices of the form
\SGARefTarget{sga3:xiii:remark:2:9:display:matrix}
\[
  g=
  \begin{pmatrix}
    t & a & c\\
    0 & 1 & b\\
    0 & 0 & 1
  \end{pmatrix},
\]
and the subgroup \(H\) formed by \SGARefLink{sga3:xiii:remark:2:9:display:matrix}{the matrices above} for which
\(b=c=0\). The Cartan subgroup of \(G\) here consists of the matrices
\(g\) for which \(a=c=0\).

\medskip
\noindent\SGARefTarget{sga3:xiii:remark:2:10}\textbf{Remarks 2.10.} ---
Let \(G\) be a smooth algebraic group over \(k\), and let \(H\) be a
smooth algebraic subgroup; we no longer assume that \(G\) and \(H\)
are connected. Suppose that \(H^0\) contains a Cartan subgroup of
\(G^0\). Then
\[
  (\mathfrak g/\mathfrak h)^H
  \subset(\mathfrak g/\mathfrak h)^{H^0}=0,
\]
so \(H^0=N^0\), where \(N\) is the normalizer of \(H\); in
particular, \(N\) is smooth.

Nevertheless, even with \(G\) connected, examples are easily
constructed in which \(H\) has a connected component \(H_i\) such
that
\[
  (\mathfrak g/\mathfrak h)^{\operatorname{ad}(h)}\neq0
  \qquad\text{for every }h\in H_i(k).
\]
Equivalently, the morphism
\[
  G\times H_i\longrightarrow G,\qquad
  (g,h)\longmapsto\operatorname{ad}(g)\cdot h
\]
is nowhere étale, or even quasi-finite. One may take, for example,
\(H\) to be the normalizer of a maximal torus in
\(\operatorname{SL}(2)_k\).

Similarly, even if the image \(H''\) of \(H\) in the finite group
\(G''=G/G^0\) is all of \(G''\), the union of the conjugates of \(H\)
in \(G\) need not be dense. For example, take \(G\) to be the
semidirect product of \(\mathbb Z/2\mathbb Z\) with
\(G^0=\operatorname{SL}(2)_k\), on which it acts by ``symmetry,'' and
take \(H\) to be the semidirect product
\(\mathbb Z/2\mathbb Z\cdot T\), where \(T\) is a maximal torus of
\(G^0\).

On the other hand, if one does not assume in advance that \(H^0\)
contains a Cartan subgroup of \(G^0\), but assumes that the union of
the conjugates of \(H\) in \(G\) is dense, then \(H^0\) necessarily
contains a Cartan subgroup of \(G^0\). To see this, one may plainly
assume that \(G\) is connected and repeat the proof of
\SGARefLink{sga3:xiii:theorem:2:1:item:vi}{2.1(vi)} \(\Rightarrow\) \SGARefLink{sga3:xiii:theorem:2:1:item:i}{(i)}, which does not require \(H\) to be
connected.
